% Unit 087 English translation of chapter6.tex, lines 2928--3141.
% Source: Methods of Algebra, Volume 2, commit
% 9a5803ff77dd3257484cb177f851a73770a59dd3 (CC BY 4.0).
% stable-unit-id: o014.aljabr2.chapter6.exercises
% Coverage: all exercises for Chapter 6; closes Chapter 6.

% segment-id: o014.li2.chapter6.exercises.h001
\begin{Exercises}
	% segment-id: o014.li2.chapter6.exercises.ex001
	\item Let $\Bbbk$ be a commutative ring and $t\in\Bbbk$. For every
	$\Bbbk$-module, write $m_t$ for the endomorphism given by multiplication
	by $t$. Let $M$ be a $G$-module. Show that the endomorphisms induced by
	$m_t\in\End_G(M)$ on $\Hm^n(G,M)$ and $\Hm_n(G,M)$ are again precisely
	$m_t$.

	\begin{hint}
		For cohomology, either write down the induced endomorphism on the
		standard complex, or take an injective resolution
		$0\to M\to I^0\to\cdots$ and show that the induced endomorphism comes
		from the diagram
		\[\begin{tikzcd}
			0 \arrow[r] & M \arrow[r] \arrow[d,"{m_t}"] &
			I^0 \arrow[r] \arrow[d,"{m_t}"] &
			I^1 \arrow[r] \arrow[d,"{m_t}"] & \cdots\\
			0 \arrow[r] & M \arrow[r] & I^0 \arrow[r] & I^1 \arrow[r] & \cdots.
		\end{tikzcd}\]
	\end{hint}

	% segment-id: o014.li2.chapter6.exercises.ex002
	\item Let the $G$-module $M$ be projective as a $\Bbbk$-module. Prove
	that $\mathsf{L}\otimes M\to\Bbbk\otimes M\simeq M$ gives a projective
	resolution of the $G$-module $M$; the same statement holds with
	$\overline{\mathsf{L}}$ in place of $\mathsf{L}$.

	\begin{hint}
		Apply Proposition~\ref{prop:group-tensor-proj}.
	\end{hint}

	% segment-id: o014.li2.chapter6.exercises.ex003
	\item Prove Proposition~\ref{prop:group-UCT} using the universal
	coefficient theorem for the cohomology of chain complexes.

	% Another candidate: Rotman, Homological Algebra. Duality Theorem 9.73. Think about it.

	% segment-id: o014.li2.chapter6.exercises.ex004
	\item Let $p$ be a prime, $G$ a $p$-group, and $M$ a $G$-module
	satisfying $pM=0$. Prove that the following properties are equivalent:
	\begin{inparaenum}[(i)]
		\item $M=0$,
		\item $M^G=0$,
		\item $M_G=0$.
	\end{inparaenum}

	\begin{hint}
		To prove (ii) $\implies$ (i), take
		$x\in M\smallsetminus\{0\}$ and consider the $\F_p$-vector subspace
		$E$ of $M$ generated by $\{gx:g\in G\}$. Use
		$|E^G|\equiv|E|\pmod p$ to show that $M^G\neq0$.

		To prove (iii) $\implies$ (i), take
		$M^\vee:=\Hom_{\F_p}(M,\F_p)$ with its natural $G$-module structure.
		Observe that
		$(M^\vee)^G\simeq\Hom_{\F_p}(M_G,\F_p)$, and deduce
		$M_G=0\implies M^\vee=0\implies M=0$.
	\end{hint}

	% segment-id: o014.li2.chapter6.exercises.ex005
	\item Let $F$ be a field. A \emph{representation} of a group $G$ is data
	$(\rho,V)$, where $V$ is an $F$-vector space, $\GL(V)$ is its group of
	linear automorphisms, and $\rho:G\to\GL(V)$ is a group homomorphism.
	Clearly, this is the same as a $G$-module with coefficients in $F$. A
	\emph{projective representation} is data $(\pi,V)$ with a homomorphism
	$\pi:G\to\PGL(V):=\GL(V)/F^\times\identity$. Isomorphisms of
	representations or projective representations are defined in the evident
	way. Henceforth identify $F^\times$ with $F^\times\identity$.
	\index{representation}
	\index{projective representation}

	\begin{enumerate}[(i)]
		\item For a projective representation $(\pi,V)$ of $G$, suppose that
		for every $g\in G$ an inverse image
		$\widetilde\pi(g)\in\GL(V)$ of $\pi(g)$ has been chosen. Verify that
		\[ \widetilde\pi(g_1)\widetilde\pi(g_2)
		=c(g_1,g_2)\widetilde\pi(g_1g_2) \]
		defines $c\in Z^2(G,F^\times)$, where $G$ acts trivially on
		$F^\times$. If $\widetilde\pi(1_G)=\identity_V$, then
		$c\in\overline{Z}^2(G,F^\times)$.
		\item Prove that the image of $c$ in $\Hm^2(G,F^\times)$ depends only
		on the isomorphism class of $(\pi,V)$ and not on any of the choices.
		\item Consider the central extension of groups
		$1\to F^\times\to\widetilde G\to G\to1$
		(Definition~\ref{def:central-ext}), where
		\[ \widetilde G:=\{(g,\widetilde\pi)\in G\times\GL(V):
		\pi(g)=\widetilde\pi\bmod F^\times\}. \]
		The first homomorphism maps $t\in F^\times$ to $(1_G,t)$ and the
		second is projection. Show that this extension is isomorphic to the
		central extension determined by the cohomology class in~(ii).
		\item Give a bijection
		\[ \{\text{representations }(\sigma,V)\text{ lifting }(\pi,V)\}
		\xleftrightarrow{1:1}
		\{\text{splittings of the central extension }\widetilde G\}. \]
		Show also that $(\pi,V)$ lifts to a representation if and only if the
		cohomology class in~(ii) is trivial.
		\item Explicitly give a group $G$ and a projective representation
		$(\pi,V)$ that cannot be lifted to a representation. Try to find an
		example with nontrivial $\pi:G\to\PGL(V)$.
	\end{enumerate}

	% segment-id: o014.li2.chapter6.exercises.ex006
	\item For subgroups $G_1\subset G_2\subset G_3$, prove the canonical
	isomorphisms of functors
	$\Ind^{G_3}_{G_2}\Ind^{G_2}_{G_1}\simeq\Ind^{G_3}_{G_1}$ and
	$\iInd^{G_3}_{G_2}\iInd^{G_2}_{G_1}\simeq\iInd^{G_3}_{G_1}$.
	Describe these isomorphisms as explicitly as possible.

	% segment-id: o014.li2.chapter6.exercises.ex007
	\item Prove that if $H$ is a subgroup of $G$, then
	$\mathrm{cd}(H)\leq\mathrm{cd}(G)$ and
	$\mathrm{hd}(H)\leq\mathrm{hd}(G)$; see
	Definition~\ref{def:group-dim}.
	\begin{hint}
		Apply Theorem~\ref{prop:Shapiro}.
	\end{hint}

	% segment-id: o014.li2.chapter6.exercises.ex008
	\item Show that for $m\in\Z_{\geq1}$,
	$\mathrm{cd}(C_m)=\mathrm{hd}(C_m)=\infty$. Prove that
	$\mathrm{cd}(G)<\infty$ implies that $G$ is torsion-free, and that
	$\mathrm{cd}(G)=0$ if and only if $G$ is the trivial group.
	\begin{hint}
		For the first part, apply
		Proposition~\ref{prop:group-cohomology-cyclic} with $A=\Z/m\Z$.
		For the other parts, consider the cyclic subgroups of $G$.
	\end{hint}

	% segment-id: o014.li2.chapter6.exercises.ex009
	\item Try to prove Corollary~\ref{prop:group-coh-torsion} directly by
	operations on the standard complex.

	% Reference: Cohomology of Number Fields, (1.5.6) and (1.5.7)
	% segment-id: o014.li2.chapter6.exercises.ex010
	\item Let $U,V$ be subgroups of $G$ and suppose that $(G:V)$ is finite.
	Write the cohomological restriction map as $\mathrm{res}^G_U$ and the
	corestriction map as $\mathrm{cor}^V_G$, and so on; the cohomological
	degree $n$ is omitted from the notation. Choose a double-coset
	decomposition
	\[ G=\bigsqcup_{\sigma\in S}U\sigma V, \]
	where the finite set $S\subset G$ consists of representatives of the
	double cosets. Let $M$ be a $G$-module.
	\begin{enumerate}[(i)]
		\item For every $\sigma\in G$ and $n\in\Z_{\geq0}$, define in a
		natural way
		\[ \mathcal{T}_\sigma:
		\Hm^n(V\cap\sigma^{-1}U\sigma,M)
		\to\Hm^n(U\cap\sigma V\sigma^{-1},M) \]
		so that for $n=0$ it becomes
		\[ M^{V\cap\sigma^{-1}U\sigma}\rightiso
		M^{U\cap\sigma V\sigma^{-1}},\qquad x\mapsto\sigma x. \]
		\item Prove that
		\[ \mathrm{res}^G_U\circ\mathrm{cor}^V_G
		=\sum_{\sigma\in S}
		\mathrm{cor}^{U\cap\sigma V\sigma^{-1}}_U
		\circ\mathcal{T}_\sigma
		\circ\mathrm{res}^V_{V\cap\sigma^{-1}U\sigma}. \]
		\item Describe the action of
		$\mathrm{res}^G_U\circ\mathrm{cor}^U_G$ when $U=V\lhd G$.
	\end{enumerate}

	% segment-id: o014.li2.chapter6.exercises.ex011
	\item Let $G$ be an abelian group, regarded as a $\Z$-module. Define
	$G\wedge G$ as the quotient of $G\dotimes{\Z}G$ by the submodule
	generated by $\{g\otimes g:g\in G\}$, and write $x\wedge y$ for the
	image of $x\otimes y$. Prove that $\Hm_2(G,\Z)\simeq G\wedge G$.

	\begin{hint}
		\index[sym1]{[,]}
		Once a presentation $G=F/R$ of the group has been fixed, the Hopf
		formula (Corollary~\ref{prop:Hopf-H2}) gives
		$\Hm_2(G,\Z)\simeq[F,F]/[F,R]$. Since $G$ is abelian,
		$[F,F]\subset R$. From the identity
		$[xy,z]=x[y,z]x^{-1}\cdot[x,z]$, the map
		\[ \psi_0:G\times G\to[F,F]/[F,R],\qquad
		(f_1R,f_2R)\mapsto[f_1,f_2]\bmod[F,R] \]
		is $\Z$-bilinear and induces a surjective homomorphism
		$\psi:G\wedge G\to[F,F]/[F,R]$. Choose a basis
		$\{x_i\}_{i\in I}$ of $F$ and give $I$ a total order. We need the
		following group-theoretic fact:
		$H:=[F,F]/[[F,F],F]$ is a free abelian group with basis the cosets of
		$[x_i,x_j]$ for $i<j$. We may therefore define
		$\phi_0:H\to G\wedge G$ by mapping the coset of $[x_i,x_j]$ to
		$x_iR\wedge x_jR$. Show that
		\[ f_1,f_2\in F\implies
		\phi_0(\text{the coset of }[f_1,f_2])=f_1R\wedge f_2R. \]
		From this, descend $\phi_0$ to a homomorphism
		$\phi:[F,F]/[F,R]\to G\wedge G$, and then show that $\phi$ and
		$\psi$ are inverse to one another.

		The group-theoretic fact above for finite $I$ appears in
		\cite[Theorem 11.2.4]{Ha76}; the general case follows from it.
	\end{hint}

	% segment-id: o014.li2.chapter6.exercises.ex012
	\item Let $H\lhd G$ and let $M$ be a $G$-module. Give the standard
	complex $C(H,M^H)$ its natural $G/H$-action, inducing a canonical
	$G/H$-action on every $\Hm^n(H,M^H)$ (see the discussion preceding
	Lemma~\ref{prop:LHS-action-aux}).

	% segment-id: o014.li2.chapter6.exercises.ex013
	\item For the filtration in Remark~\ref{rem:LHS-SS-mult}, describe the
	$E_2$ page of its spectral sequence and show that the properties in
	Theorem~\ref{prop:LHS-SS} continue to hold. Explain the compatibility of
	the spectral sequence with the cup product as fully as possible.

	% segment-id: o014.li2.chapter6.exercises.ex014
	\item Consider the finite cyclic group $C_m$ of order $m$. For every
	$C_m$-module $A$ and $n\in\Z_{\geq1}$, prove that the periodicity
	isomorphism $\Hm^n(C_m,A)\to\Hm^{n+2}(C_m,A)$ in
	Proposition~\ref{prop:group-cohomology-cyclic} equals the composite of
	the two connecting homomorphisms from the canonical short exact sequences
	in Example~\ref{eg:cup-periodicity}:
	\begin{gather*}
		0\to\mathfrak{I}\otimes A\to\Z[C_m]\otimes A\to A\to0,\\
		0\to A\xrightarrow{\nu\otimes\identity}\Z[C_m]\otimes A
		\xrightarrow{(\sigma-1)\otimes\identity}\mathfrak{I}\otimes A\to0.
	\end{gather*}

	\begin{hint}
		Compute cohomology with the periodic complex
		$A\xrightarrow{\sigma-1}A\xrightarrow{\nu}A\to\cdots$ from the proof
		of Proposition~\ref{prop:group-cohomology-cyclic}, starting in degree
		zero, and use it to describe the connecting homomorphisms. Treat the odd
		and even cases separately; some computation is required.
	\end{hint}

	% segment-id: o014.li2.chapter6.exercises.ex015
	\item Let $\Bbbk$ be any commutative ring and $m\in\Z_{\geq1}$.
	Determine the structure of the cohomology algebra
	$\Hm^\bullet(C_m,\Bbbk)$ as a graded $\Bbbk$-algebra.
%	\begin{hint}
%		Substituting $A=\Bbbk$ in Example~\ref{eg:cup-periodicity} may help.
%	\end{hint}

	% Reference: Spanier, p.254
	% According to May (p.155), it is rather $\lrangle{\alpha \cup \beta, \gamma} = \lrangle{\beta, \alpha \cap \gamma}$
	% segment-id: o014.li2.chapter6.exercises.ex016
	\item For the cap product introduced in Remark~\ref{rem:cap-product},
	prove the following adjunction relation with the cup product:
	% Editorial correction: the frozen Chinese source and Indonesian witness use
	% an unbound \gamma on the left; the homology class introduced below is x.
	\[ \lrangle{\alpha\cup\beta,x}
	=\lrangle{\alpha,\beta\cap x}\in(M_1\otimes M_2\otimes M_3)_G, \]
	where $M_1,M_2,M_3$ are $G$-modules,
	$\alpha\in\Hm^p(G,M_1)$, $\beta\in\Hm^q(G,M_2)$, and
	$x\in\Hm_{p+q}(G,M_3)$.

	% segment-id: o014.li2.chapter6.exercises.ex017
	\item In this problem $G$ is a finite group, the coefficient ring in
	cohomology and homology is $\Z$, and a free $G$-module means a free
	$\Z[G]$-module. Write the contragredient module of a $G$-module $M$ as
	$M^*:=\Hom_{\Z}(M,\Z)$.
	\begin{enumerate}[(i)]
		\item Take a free resolution $\epsilon:P\to\Z$ of the trivial
		$G$-module, where $P:=[\cdots\to P_n\to P_{n-1}\to\cdots]$ and every
		$P_n$ is finitely generated. Define
		$P^*:=[\cdots\to P_n^*\to P_{n+1}^*\to\cdots]$. Show that
		$\epsilon^*:\Z\to P^*$ is a quasi-isomorphism.
		\begin{hint}
			Every $P_n$ is a free $\Z$-module. If the chain complex of free
			$\Z$-modules $(C_n,\partial_n)_{n\in\Z}$ is exact and
			$Z_n:=\Ker(\partial_n)$, then the short exact sequence
			$0\to Z_{n+1}\to C_{n+1}\xrightarrow{\partial_n}Z_n\to0$
			splits for every $n$.
		\end{hint}
		\item Prove that every $P_n^*$ is a finitely generated projective
		$G$-module.
		\begin{hint}
			It is enough to show that $\Z[G]^*\simeq\Z[G]$.
		\end{hint}
		\item Let $Q$ be a finitely generated projective $G$-module. For
		every $G$-module $M$, prove that the canonical morphism
		$\nu:(Q\otimes M)_G\to(Q\otimes M)^G$ from
		\eqref{eqn:nu-coinv-inv} is an isomorphism.
		\begin{hint}
			Use a direct-sum decomposition to reduce to the case $Q=\Z[G]$, and
			then verify it directly. For greater simplicity, the automorphism
			\[ \left(\sum_{g\in G}a_gg\right)\otimes x
			\mapsto\sum_{g\in G}a_g(g\otimes g^{-1}x) \]
			turns the diagonal $G$-action on $\Z[G]\otimes M$ into left
			multiplication.
		\end{hint}

		\item Define $P_{-n}:=P^*_{n-1}$ and splice $P$ and $P^*$ into the
		chain complex
		\[ \mathcal{P}:=[\cdots\to P_1\to P_0
		\xrightarrow{\epsilon^*\epsilon}P_{-1}\to P_{-2}\to\cdots]. \]
		Prove that for every $G$-module $M$ there is a canonical isomorphism
		\[ \Hm^n\Hom_G(\mathcal{P},M)
		\simeq\text{Tate cohomology }\TaHm^n(G,M),\qquad n\in\Z. \]

		\begin{hint}
			For $n\geq1$ the statement is clear. For $n\leq-2$, use the
			following property: if $Q$ is a finitely generated projective
			$G$-module, then $Q\otimes M\rightiso\Hom(Q^*,M)$. Deduce that
			\[ (Q\otimes M)_G\xrightarrow[\nu]{\sim}(Q\otimes M)^G
			\rightiso\Hom(Q^*,M)^G=\Hom_G(Q^*,M). \]
			For $n=0,-1$, identify
			$\Hom_G(P_{-1},M)\to\Hom_G(P_0,M)$ by means of the isomorphism
			above with
			\[ (P_0\otimes M)_G\xrightarrow{(\epsilon\otimes\identity)_G}M_G
			\xrightarrow{\nu}M^G\xrightarrow{\epsilon^*}\Hom_G(P_0,M). \]
		\end{hint}

		\item Use this result to prove the long exact sequence in Tate
		cohomology directly.
		\item Take the free resolution $\epsilon:P\to\Z$ from
		Proposition~\ref{prop:trivial-mod-nor-resolution}, choose its natural
		basis for every $P_n$, and then explicitly describe the chain complex
		$\mathcal{P}$ from~(iv).
	\end{enumerate}

	Based on~(vi), a cup product on Tate cohomology can be defined as in
	Lemma~\ref{prop:cup-resolution}, by appropriately defining
	$\Delta:\mathcal{P}\to\mathcal{P}\otimes\mathcal{P}$. For explicit
	formulas, see \cite[Chapter IV, \S 7]{CF67}.
	\index{cup product}
	\index{Tate cohomology}
	\index[sym1]{HnGMTate}

	% segment-id: o014.li2.chapter6.exercises.ex018
	\item Let $G$ be a finite group and $H$ a subgroup. For a $G$-module
	$M$, use dimension shifting (Corollary~\ref{prop:Tate-dim-shift}) to
	extend the canonical morphisms from \S\ref{sec:group-finite-index},
	$\mathrm{res}^n:\TaHm^n(G,M)\to\TaHm^n(H,M)$ and
	$\mathrm{res}_n:\TaHm^{-n-1}(H,M)\to\TaHm^{-n-1}(G,M)$, from the case
	$n\geq1$ to all $n\in\Z$. For the moment, denote the former by
	$\mathrm{res}_+$ and the latter by $\mathrm{res}_-$.
	\begin{enumerate}[(i)]
		\item Describe the action of $\mathrm{res}_+$ on $\TaHm^0$ and the
		action of $\mathrm{res}_-$ on $\TaHm^{-1}$.
		\item Prove that $\mathrm{res}_+$ on $\TaHm^{-1}$ is induced by
		$\nu_{G|H}:M_G\to M_H$ (Definition~\ref{def:cor-zero}), and on
		$\TaHm^{<-1}$ equals $\mathrm{cor}_n$
		(Definition--Proposition~\ref{def:cor}).
		\item Prove that $\mathrm{res}_-$ on $\TaHm^0$ is induced by
		$\nu^{G|H}:M^H\to M^G$, and on $\TaHm^{>0}$ equals
		$\mathrm{cor}^n$.
	\end{enumerate}

	% segment-id: o014.li2.chapter6.exercises.ex019
	\item For a group $G$ and a bounded-below complex of $G$-modules
	$M=[\cdots\to M^p\to M^{p+1}\to\cdots]$, let $p$ vary in the standard
	complexes $C(G,M^p)$ to form a double complex. Denote its total complex by
	$C(G,M):=\tot((C^q(G,M^p))_{p,q})$. Prove the canonical isomorphism
	\[ \Hm^n(C(G,M))\simeq
	\Hm^n(G,[\cdots\to M^p\to M^{p+1}\to\cdots]),\qquad n\in\Z, \]
	where the right-hand side is defined as in \S\ref{sec:derived-primer} and
	is also called group hypercohomology.

	\begin{hint}
		By definition, the right-hand side is obtained by taking an injective
		resolution $M\to I$ (recall that $I$ is a bounded-below complex, every
		$I^q$ is an injective $G$-module, and $M\to I$ is a quasi-isomorphism)
		and then taking the cohomology determined by the $\Hom$ complex
		$\Hom^\bullet_G(\Bbbk,I)$. Show that the same cohomology can be
		computed by $\Hom^\bullet_G(\mathsf{L},M)$.
	\end{hint}

	% segment-id: o014.li2.chapter6.exercises.ex020
	\item Prove the corresponding result when $G$ is a profinite group,
	requiring every $M^p$ to be a smooth $G$-module and using the continuous
	standard complexes $C(G,M^p)$.

	\begin{hint}
		Define a functor $C(G,\cdot)$ from
		$\cate{C}^+(G\dcate{Mod}^\infty)$ to
		$\cate{C}^+(\Bbbk\dcate{Mod})$. Show that it induces a triangulated
		functor at the level of $\cate{K}^+(\cdots)$ and preserves acyclic
		objects, and hence induces a triangulated functor at the level of
		$\cate{D}^+(\cdots)$, still denoted by $C(G,\cdot)$. The natural
		transformation $(\cdot)^G\to C(G,\cdot)$ and the universal property of
		right derived functors give a morphism of triangulated functors
		$\mathrm{R}(\cdot)^G\to C(G,\cdot)$. To prove that this is an
		isomorphism, use the way-out lemma (Proposition~\ref{prop:wayout}) to
		reduce to the case handled by
		Theorem~\ref{prop:profinite-group-standard}.
	\end{hint}

	% segment-id: o014.li2.chapter6.exercises.ex021
	\item For a profinite group $G$ and a finite $G$-module $M$, give a
	complete proof of the bijective correspondence mentioned in
	\S\ref{sec:profinite-cohomology} between equivalence classes of
	profinite group extensions and $\Hm^2(G,M)$.

	\begin{hint}
		Follow the argument in \S\ref{sec:grp-low-degree}. In the profinite
		case the required observations are: (i) the group structure defined by
		every $f\in\overline{Z}^2(G,M)$ on the product space $M\times G$ makes
		it a profinite group; (ii) for a profinite group extension $E$, there
		is an isomorphism of topological groups $E/M\rightiso G$; (iii) in
		this situation, Lemma~\ref{prop:profinite-section} gives a continuous
		% Editorial correction: the frozen Chinese source and Indonesian witness
		% repeat the quotient map's type; its section has type G -> E.
		section $s:G\to E$ that can be adjusted to be
		normalized.
	\end{hint}

	% segment-id: o014.li2.chapter6.exercises.ex022
	\item For any field $F$, define $\GL(n,F)$ to be the group of invertible
	$n\times n$ matrices over $F$, and $\SL(n,F)$ to be the subgroup cut out
	by $\det=1$. Observe that $G:=\Gal(\CC|\R)$ acts on $\GL(n,\CC)$ and
	$\SL(n,\CC)$, with fixed points $\GL(n,\R)$ and $\SL(n,\R)$,
	respectively.
	\begin{enumerate}[(i)]
		\item Take $n=2$. Let $B:=\SL(2,\CC)$ act by conjugation on the space
		of all $2\times2$ matrices, and define
		$A:=\Stab(\begin{smallmatrix}0&-1\\1&0\end{smallmatrix})$. Show that
		this subgroup is invariant under the $G$-action and describe $A^G$.
		\item Show that $B/A$ (respectively, $B^G/A^G$) is naturally
		identified with the orbit of
		$\left(\begin{smallmatrix}0&-1\\1&0\end{smallmatrix}\right)$ under
		conjugation by $\SL(2,\CC)$ (respectively, $\SL(2,\R)$). For $B/A$,
		explain how the $G$-action is reflected on the orbit.
		\item Show that $(B/A)^G\neq B^G/A^G$. In other words, there is a
		real $2\times2$ matrix that is conjugate to
		$\left(\begin{smallmatrix}0&-1\\1&0\end{smallmatrix}\right)$ by
		$\SL(2,\CC)$ but not by $\SL(2,\R)$.
		\item Show that the phenomenon above does not occur if $\SL$ is
		replaced by $\GL$.
	\end{enumerate}

	% segment-id: o014.li2.chapter6.exercises.ex023
	\item State and prove the compatibility between the isomorphisms in
	Theorem~\ref{prop:Shapiro-noncommutative} and the connecting
	homomorphisms in Theorem~\ref{prop:nonabelian-long}.

	% segment-id: o014.li2.chapter6.exercises.ex024
	\item (``Twisting'' group cohomology) Use the notation of
	\S\ref{sec:nonabelian-coh}. Let $A$ be a $G$-group.
	\begin{enumerate}[(i)]
		\item For $c\in Z^1(G,A)$, define a new group action $G\times A\to A$
		by
		\[ {}^{c,g}a:=c(g)\cdot{}^ga\cdot c(g)^{-1},
		\qquad g\in G,\ a\in A. \]
		Show that this indeed gives a $G$-group, denoted by ${}^cA$. Prove
		also that if $c\sim c'$, there is an isomorphism of $G$-groups
		${}^{c'}A\simeq{}^cA$; describe it explicitly.
		\item Prove that every $c\in Z^1(G,A)$ induces a bijection
		\begin{align*}
			\nu_c:Z^1(G,{}^cA)&\xrightarrow{1:1}Z^1(G,A),\\
			c'&\longmapsto c'c\quad\text{(pointwise multiplication)},
		\end{align*}
		which induces a bijection
		${\bm\nu}_c:\Hm^1(G,{}^cA)\xrightarrow{1:1}\Hm^1(G,A)$ that maps the
		basepoint to $[c]$. Give a simplified description when $A$ is abelian.
		\item Suppose there is a short exact sequence
		$1\to A\to B\to B/A\to1$ as in
		Theorem~\ref{prop:nonabelian-long}~(i). For $c\in Z^1(G,A)$, prove
		that the following diagram commutes:
		\[\begin{tikzcd}
			\Hm^1(G,{}^cA) \arrow[r] \arrow[d,"{{\bm\nu}_c}"'] &
			\Hm^1(G,{}^cB) \arrow[d,"{\bm{\nu}_c}"]\\
			\Hm^1(G,A) \arrow[r] & \Hm^1(G,B).
		\end{tikzcd}\]
		Use this to describe all fibers of the map
		$\Hm^1(G,A)\to\Hm^1(G,B)$.
		\item State the modifications required when $G$ is a profinite group.
	\end{enumerate}
	\index{group cohomology!twist}
\end{Exercises}
